\paragraph{Theorem F.2.22 (Proved: a terminating lattice algorithm for each rational pair).}
Let $(A,C)$ be any irreducible rational rank-five pair in
one of the eight spectral cases of F.2.14, with
$A^5=I$, $B=C^{-1}$, and $(AB-I)^2=0$ of rank two.
There is an explicit terminating algorithm which decides
whether the pair preserves an integral lattice, and, if so,
lists all such lattices up to simultaneous integral conjugacy.
It then lists all rank-six extensions with the same finite
orders and primitive maximal square-zero cusp, up to
simultaneous integral conjugacy.  The bounds below depend
on the input rational pair and are proved to lose no lattice.
They do not bound the rational parameter space or classify
all affine twists and resulting homology.

\begin{proof}
We give the algorithm and its termination proof.
The fixed line of $A$ has dimension one.  Choose a nonzero
vector $e$ on it and a nonzero fixed covector $\varphi$.
Irreducibility implies that the vectors $ge$, with $g$ a
word in $A,C$, span the whole rational space; the analogous
statement holds for the covectors $\varphi h$.
In fact words of length at most four suffice.  The spans of
words of lengths at most $j$ increase from dimension one;
if one step does not increase the span, it is invariant
under both generators and hence already the whole space.
There are only four possible increases.

Choose words $g_1,\ldots,g_5$ and $h_1,\ldots,h_5$ giving
bases in this way, and put
\[
 P=I+A+\cdots+A^4,\qquad T_{ij}=g_iPh_j.
 \tag{F.L1}
\]
The rank-one operator $P$ is a nonzero scalar multiple of
$e\varphi$.  Hence the 25 matrices $T_{ij}$ are a rational
basis of $\operatorname{Mat}_5(\mathbb Q)$.
Each belongs to $\mathcal R=\mathbb Z[A,C]$.
Let $\mathcal R_0$ be their integral span and set
\[
 \mathcal R_0^\#=
 \{X:\operatorname{tr}(XT)\in\mathbb Z
                     \text{ for all }T\in\mathcal R_0\}.
 \tag{F.L2}
\]
The trace pairing on the full matrix algebra is
nondegenerate, as is seen from the matrix units, so this
is another full rational lattice, computed by inverting
the 25-by-25 Gram matrix.

If any stable lattice exists, every element of $\mathcal R$
acts integrally on it.  Consequently
$\operatorname{tr}(XY)\in\mathbb Z$ for all
$X,Y\in\mathcal R$, and necessarily
\[
 \mathcal R_0\subset\mathcal R\subset\mathcal R_0^\#.
 \tag{F.L3}
\]
A nonintegral Gram entry therefore proves nonexistence.
Otherwise let
$d=[\mathcal R_0^\#:\mathcal R_0]$, the absolute Gram
determinant.  Start with the lattice generated by
$\mathcal R_0$ and $I$, and repeatedly replace a current
lattice $L$ by $L+AL+CL$, using exact Hermite reduction.
If it leaves $\mathcal R_0^\#$, (F.L3) proves that there
is no stable lattice.
While it remains inside, every strict enlargement reduces
its index by a factor at least two.
There can be at most $\lfloor\log_2d\rfloor$ such
enlargements starting from $\mathcal R_0$, so a bound of
$1+\lfloor\log_2d\rfloor$ tests after adjoining $I$ is
sufficient.  Stabilization gives exactly $\mathcal R$:
the stable lattice contains $I$ and all words in $A,C$,
and every element introduced belongs to their ring.
It is then an order, and
$\Lambda_0=\mathcal R\mathbb Z^5$ is an explicitly
computed stable lattice.  This proves the decision step,
including termination in the nonexistence case.

Use a basis of $\Lambda_0$ from now on, so
$\mathcal R\subset\operatorname{Mat}_5(\mathbb Z)$.
Smith reduction gives an integer $c\ge1$ with
\[
 c\operatorname{Mat}_5(\mathbb Z)\subset\mathcal R;
 \tag{F.L4}
\]
one may take the largest Smith factor of the inclusion
of the order in the full matrix lattice.
Normalize any stable lattice by a rational homothety so
that it is contained in $\mathbb Z^5$ and the gcd of all
coordinates of a lattice basis is one.
The operators $cE_{ij}$ preserve it.  The coordinate values
of its vectors generate $\mathbb Z$, so Bezout's identity,
applied to $cE_{ij}$, shows that it contains every $ce_i$.
Therefore every normalized lattice lies in the finite interval
\[
 c\mathbb Z^5\subset\Lambda\subset\mathbb Z^5.
 \tag{F.L5}
\]
Conversely every lattice in this interval stable under
$A,C$ is retained by the algorithm.
The only rational operators commuting with the pair are
scalars: they act by a scalar on the one-dimensional fixed
line of $A$, and then by the same scalar on the spanning
vectors $g_ie$.  Thus rational equivalence of two stable
lattices is precisely homothety.  The coordinate-gcd
normalization is unique, so (F.L5) gives the integral
conjugacy classification without further identifications.

Here is a finite enumeration bound which is also implemented.
For each $p^a\Vert c$, enumerate all subspaces of
$\mathbb F_p^5$ in reduced echelon form.
Beginning with $\mathbb Z^5$, pass through $a$ levels of
stable lattices between $pH$ and $H$, retaining only those
containing $p^a\mathbb Z^5$ and identifying equal Hermite
forms.  Every desired lattice occurs: for any target
$\Lambda$, the chain
$\Lambda+p^j\mathbb Z^5$, $0\le j\le a$, is such a
chain and ends at $\Lambda$.
If $s_p$ denotes the number of all subspaces of
$\mathbb F_p^5$, there are at most $s_p^a$ branches before
deduplication.  Intersect the resulting lattices for the
different primes and discard those of coordinate gcd
greater than one.  The Chinese remainder theorem proves
that this enumerates (F.L5) completely.

For each quotient lattice $W_{\mathbb Z}$ so obtained,
put $N_W=AB-I$.
Its image must be primitive: for a primitive full cusp,
$\operatorname{im}N_V=\ker N_V$ contains the primitive
invariant line, and quotienting that image by the line
gives a primitive image in $W_{\mathbb Z}$.
Reject a quotient whose two nonzero Smith factors are
not one.  For the remaining quotients use extension rows
$c_1,c_2\in\mathbb Z^{1\times5}$ and impose
\[
\begin{gathered}
 c_1(I+A+\cdots+A^4)=0,
 \qquad c_2(I+B+\cdots+B^{m-1})=0,\\
 (c_2+c_1B)N_W=0,
 \qquad m=\operatorname{ord}(B).
\end{gathered}
 \tag{F.L6}
\]
These are integral linear equations.  Their integral kernel
modulo the coboundaries
\[
 z\longmapsto\bigl(z(A-I),z(B-I)\bigr),
 \qquad z\in\mathbb Z^{1\times5},
 \tag{F.L7}
\]
is a finitely generated group $G$ of rank one by F.2.11.
Compute its free and torsion factors by Smith reduction.
Choose an integral generator $w$ of the free rank-one
quotient $\ker N_W/\operatorname{im}N_W$.
The homomorphism
\[
 \kappa:G\longrightarrow\mathbb Z,
 \qquad [(c_1,c_2)]\longmapsto(c_2+c_1B)w
 \tag{F.L8}
\]
is well defined and nonzero: (F.L6) kills the image, and
(F.L7) kills $w$; rational nonvanishing is F.2.11.
The full cusp has unit Smith factors exactly when
$\kappa=\pm1$.
Indeed the quotient rank-two image is already primitive,
and the remaining invariant-line image is generated by
that one integer.
If $\operatorname{im}\kappa\ne\mathbb Z$, reject the
lattice.  Otherwise the inverse image of $1$ is a finite
coset of $\operatorname{Tor}G$, and lists all extensions.
Changing the sign of the invariant-line generator identifies
$\kappa=-1$ with $\kappa=1$.
There are no further identifications: an integral isomorphism
on the quotient is $\pm I$, on the invariant line is
$\pm1$, and, after $\kappa=1$ is fixed, their ratio is one.
The remaining row is exactly an integral change of splitting
(F.L7).
This proves completeness and termination of the extension
step as well.

The order, trace-dual bound, conductor and finite-interval
steps are implemented in
\texttt{fourfolds\_lattice\_algorithm.py};
the extension-kernel and Smith steps are implemented in
\texttt{fourfolds\_two\_classes\_lattices.py} with the
finite order as an argument.  Their applications and exact
exhaustive bounds are recorded next.
\end{proof}

\paragraph{Theorem F.2.23 (Proved: complete lattice classes for the two family survivors).}
For each of the two rational representations $v=0$ and $v=1$
of F.2.16, there are exactly two homothety classes of stable
rank-five quotient lattices, and exactly three integral
conjugacy classes of rank-six extensions with primitive
maximal cusp.  Exactly one of those three can have simply
connected finite affine fillings.  It is the integral class
already realized in F.2.6, respectively F.2.18.
Thus arbitrary stable-lattice changes produce no further
linear survivor classes anywhere in the family F.2.16.
This assertion retains the classification of all affine
twists and multiplicities as a separate problem.

\begin{proof}
Apply the preceding algorithm in the original five-cycle
lattice (F.U1).  For both parameters the order is obtained
after four strict enlargements.  Its Smith factors in the
full matrix lattice are 21 factors one and four factors two;
thus the conductor $c=2$ in (F.L4) is exact.
All 374 subspaces of $\mathbb F_2^5$, including zero and
the whole space, are enumerated, with no height restriction
or neighbor-depth assumption.
After imposing stability and coordinate gcd one, the two
lattice representatives are
\[
 \mathbb Z^5,\qquad
 \Lambda_* =2\mathbb Z^5+\mathbb Z(1,1,1,1,1)^t.
 \tag{F.L9}
\]
The latter has index $16$.
The generic prime-by-prime implementation of (F.L5) and
an independent bit-vector enumeration give exactly the
same two Hermite forms.

For $\mathbb Z^5$, the extension group in (F.L6)--(F.L7)
is $\mathbb Z$ and $\kappa$ is a unit on its generator.
For $\Lambda_*$, that group is
$\mathbb Z\oplus\mathbb Z/2$, again with primitive
$\kappa$.  Hence there are exactly one and two primitive
extensions, respectively.
For each extension the program computes the image of the
primitive global invariant homology covector $\ell$ on
the two finite fixed lattices.  The resulting positive
divisibilities are
\[
\begin{array}{c|c|c}
 \text{quotient lattice}&\text{extension class}&
       (\operatorname{div}_1,\operatorname{div}_2)\\\hline
 \mathbb Z^5&\text{unique}&(1,2)\\
 \Lambda_*&\text{first torsion coset}&(1,1)\\
 \Lambda_*&\text{second torsion coset}&(1,2).
\end{array}
 \tag{F.L10}
\]
The table is identical for $v=0$ and $v=1$.
All six full matrices, the integral kernel bases, the
Smith factors and the finite divisibilities are computed
exactly by \texttt{fourfolds\_two\_classes\_lattices.py}
and saved in its output and JSON certificate.

For a simply connected filling the determinant quotient
of F.2.12 requires $\gcd(\ell(v_i),m_i)=1$ at each
finite cap.  The second affine multiplicity is divisible
by four.  A divisibility of two in (F.L10) therefore
excludes that class for every choice of twists and larger
multiplicities.
The sole remaining class has an explicit integral
intertwiner of determinant $\pm1$ with the corresponding
pair in (F.U6), as printed and checked by the program.
Those pairs already have the verified simply connected,
three-connected fillings.  F.2.19 excludes every other
rational parameter for all lattices.
Together these facts prove all the claims, including the
absence of any additional linear survivor class in the
unbounded family.
\end{proof}
